% Options for packages loaded elsewhere
\PassOptionsToPackage{unicode}{hyperref}
\PassOptionsToPackage{hyphens}{url}
\PassOptionsToPackage{dvipsnames,svgnames,x11names}{xcolor}
%
\documentclass[
  letterpaper,
  DIV=11,
  numbers=noendperiod]{scrartcl}

\usepackage{amsmath,amssymb}
\usepackage{iftex}
\ifPDFTeX
  \usepackage[T1]{fontenc}
  \usepackage[utf8]{inputenc}
  \usepackage{textcomp} % provide euro and other symbols
\else % if luatex or xetex
  \usepackage{unicode-math}
  \defaultfontfeatures{Scale=MatchLowercase}
  \defaultfontfeatures[\rmfamily]{Ligatures=TeX,Scale=1}
\fi
\usepackage{lmodern}
\ifPDFTeX\else  
    % xetex/luatex font selection
\fi
% Use upquote if available, for straight quotes in verbatim environments
\IfFileExists{upquote.sty}{\usepackage{upquote}}{}
\IfFileExists{microtype.sty}{% use microtype if available
  \usepackage[]{microtype}
  \UseMicrotypeSet[protrusion]{basicmath} % disable protrusion for tt fonts
}{}
\makeatletter
\@ifundefined{KOMAClassName}{% if non-KOMA class
  \IfFileExists{parskip.sty}{%
    \usepackage{parskip}
  }{% else
    \setlength{\parindent}{0pt}
    \setlength{\parskip}{6pt plus 2pt minus 1pt}}
}{% if KOMA class
  \KOMAoptions{parskip=half}}
\makeatother
\usepackage{xcolor}
\setlength{\emergencystretch}{3em} % prevent overfull lines
\setcounter{secnumdepth}{5}
% Make \paragraph and \subparagraph free-standing
\makeatletter
\ifx\paragraph\undefined\else
  \let\oldparagraph\paragraph
  \renewcommand{\paragraph}{
    \@ifstar
      \xxxParagraphStar
      \xxxParagraphNoStar
  }
  \newcommand{\xxxParagraphStar}[1]{\oldparagraph*{#1}\mbox{}}
  \newcommand{\xxxParagraphNoStar}[1]{\oldparagraph{#1}\mbox{}}
\fi
\ifx\subparagraph\undefined\else
  \let\oldsubparagraph\subparagraph
  \renewcommand{\subparagraph}{
    \@ifstar
      \xxxSubParagraphStar
      \xxxSubParagraphNoStar
  }
  \newcommand{\xxxSubParagraphStar}[1]{\oldsubparagraph*{#1}\mbox{}}
  \newcommand{\xxxSubParagraphNoStar}[1]{\oldsubparagraph{#1}\mbox{}}
\fi
\makeatother

\usepackage{color}
\usepackage{fancyvrb}
\newcommand{\VerbBar}{|}
\newcommand{\VERB}{\Verb[commandchars=\\\{\}]}
\DefineVerbatimEnvironment{Highlighting}{Verbatim}{commandchars=\\\{\}}
% Add ',fontsize=\small' for more characters per line
\usepackage{framed}
\definecolor{shadecolor}{RGB}{241,243,245}
\newenvironment{Shaded}{\begin{snugshade}}{\end{snugshade}}
\newcommand{\AlertTok}[1]{\textcolor[rgb]{0.68,0.00,0.00}{#1}}
\newcommand{\AnnotationTok}[1]{\textcolor[rgb]{0.37,0.37,0.37}{#1}}
\newcommand{\AttributeTok}[1]{\textcolor[rgb]{0.40,0.45,0.13}{#1}}
\newcommand{\BaseNTok}[1]{\textcolor[rgb]{0.68,0.00,0.00}{#1}}
\newcommand{\BuiltInTok}[1]{\textcolor[rgb]{0.00,0.23,0.31}{#1}}
\newcommand{\CharTok}[1]{\textcolor[rgb]{0.13,0.47,0.30}{#1}}
\newcommand{\CommentTok}[1]{\textcolor[rgb]{0.37,0.37,0.37}{#1}}
\newcommand{\CommentVarTok}[1]{\textcolor[rgb]{0.37,0.37,0.37}{\textit{#1}}}
\newcommand{\ConstantTok}[1]{\textcolor[rgb]{0.56,0.35,0.01}{#1}}
\newcommand{\ControlFlowTok}[1]{\textcolor[rgb]{0.00,0.23,0.31}{\textbf{#1}}}
\newcommand{\DataTypeTok}[1]{\textcolor[rgb]{0.68,0.00,0.00}{#1}}
\newcommand{\DecValTok}[1]{\textcolor[rgb]{0.68,0.00,0.00}{#1}}
\newcommand{\DocumentationTok}[1]{\textcolor[rgb]{0.37,0.37,0.37}{\textit{#1}}}
\newcommand{\ErrorTok}[1]{\textcolor[rgb]{0.68,0.00,0.00}{#1}}
\newcommand{\ExtensionTok}[1]{\textcolor[rgb]{0.00,0.23,0.31}{#1}}
\newcommand{\FloatTok}[1]{\textcolor[rgb]{0.68,0.00,0.00}{#1}}
\newcommand{\FunctionTok}[1]{\textcolor[rgb]{0.28,0.35,0.67}{#1}}
\newcommand{\ImportTok}[1]{\textcolor[rgb]{0.00,0.46,0.62}{#1}}
\newcommand{\InformationTok}[1]{\textcolor[rgb]{0.37,0.37,0.37}{#1}}
\newcommand{\KeywordTok}[1]{\textcolor[rgb]{0.00,0.23,0.31}{\textbf{#1}}}
\newcommand{\NormalTok}[1]{\textcolor[rgb]{0.00,0.23,0.31}{#1}}
\newcommand{\OperatorTok}[1]{\textcolor[rgb]{0.37,0.37,0.37}{#1}}
\newcommand{\OtherTok}[1]{\textcolor[rgb]{0.00,0.23,0.31}{#1}}
\newcommand{\PreprocessorTok}[1]{\textcolor[rgb]{0.68,0.00,0.00}{#1}}
\newcommand{\RegionMarkerTok}[1]{\textcolor[rgb]{0.00,0.23,0.31}{#1}}
\newcommand{\SpecialCharTok}[1]{\textcolor[rgb]{0.37,0.37,0.37}{#1}}
\newcommand{\SpecialStringTok}[1]{\textcolor[rgb]{0.13,0.47,0.30}{#1}}
\newcommand{\StringTok}[1]{\textcolor[rgb]{0.13,0.47,0.30}{#1}}
\newcommand{\VariableTok}[1]{\textcolor[rgb]{0.07,0.07,0.07}{#1}}
\newcommand{\VerbatimStringTok}[1]{\textcolor[rgb]{0.13,0.47,0.30}{#1}}
\newcommand{\WarningTok}[1]{\textcolor[rgb]{0.37,0.37,0.37}{\textit{#1}}}

\providecommand{\tightlist}{%
  \setlength{\itemsep}{0pt}\setlength{\parskip}{0pt}}\usepackage{longtable,booktabs,array}
\usepackage{calc} % for calculating minipage widths
% Correct order of tables after \paragraph or \subparagraph
\usepackage{etoolbox}
\makeatletter
\patchcmd\longtable{\par}{\if@noskipsec\mbox{}\fi\par}{}{}
\makeatother
% Allow footnotes in longtable head/foot
\IfFileExists{footnotehyper.sty}{\usepackage{footnotehyper}}{\usepackage{footnote}}
\makesavenoteenv{longtable}
\usepackage{graphicx}
\makeatletter
\newsavebox\pandoc@box
\newcommand*\pandocbounded[1]{% scales image to fit in text height/width
  \sbox\pandoc@box{#1}%
  \Gscale@div\@tempa{\textheight}{\dimexpr\ht\pandoc@box+\dp\pandoc@box\relax}%
  \Gscale@div\@tempb{\linewidth}{\wd\pandoc@box}%
  \ifdim\@tempb\p@<\@tempa\p@\let\@tempa\@tempb\fi% select the smaller of both
  \ifdim\@tempa\p@<\p@\scalebox{\@tempa}{\usebox\pandoc@box}%
  \else\usebox{\pandoc@box}%
  \fi%
}
% Set default figure placement to htbp
\def\fps@figure{htbp}
\makeatother

\usepackage{booktabs}
\usepackage{longtable}
\usepackage{array}
\usepackage{multirow}
\usepackage{wrapfig}
\usepackage{float}
\usepackage{colortbl}
\usepackage{pdflscape}
\usepackage{tabu}
\usepackage{threeparttable}
\usepackage{threeparttablex}
\usepackage[normalem]{ulem}
\usepackage{makecell}
\usepackage{xcolor}
\usepackage{xurl} % break long URLs
\PassOptionsToPackage{hyphens}{url}
\usepackage{fvextra} % better verbatim handling (no Highlighting redefinition)
\usepackage{booktabs} % professional table rules
\KOMAoption{captions}{tableheading}
\makeatletter
\@ifpackageloaded{caption}{}{\usepackage{caption}}
\AtBeginDocument{%
\ifdefined\contentsname
  \renewcommand*\contentsname{Table of contents}
\else
  \newcommand\contentsname{Table of contents}
\fi
\ifdefined\listfigurename
  \renewcommand*\listfigurename{List of Figures}
\else
  \newcommand\listfigurename{List of Figures}
\fi
\ifdefined\listtablename
  \renewcommand*\listtablename{List of Tables}
\else
  \newcommand\listtablename{List of Tables}
\fi
\ifdefined\figurename
  \renewcommand*\figurename{Figure}
\else
  \newcommand\figurename{Figure}
\fi
\ifdefined\tablename
  \renewcommand*\tablename{Table}
\else
  \newcommand\tablename{Table}
\fi
}
\@ifpackageloaded{float}{}{\usepackage{float}}
\floatstyle{ruled}
\@ifundefined{c@chapter}{\newfloat{codelisting}{h}{lop}}{\newfloat{codelisting}{h}{lop}[chapter]}
\floatname{codelisting}{Listing}
\newcommand*\listoflistings{\listof{codelisting}{List of Listings}}
\makeatother
\makeatletter
\makeatother
\makeatletter
\@ifpackageloaded{caption}{}{\usepackage{caption}}
\@ifpackageloaded{subcaption}{}{\usepackage{subcaption}}
\makeatother

\usepackage{bookmark}

\IfFileExists{xurl.sty}{\usepackage{xurl}}{} % add URL line breaks if available
\urlstyle{same} % disable monospaced font for URLs
\hypersetup{
  pdftitle={Social Network Analysis of the Canadian Disability Rights Community},
  pdfauthor={Dr.~Derrick L. Cogburn},
  colorlinks=true,
  linkcolor={blue},
  filecolor={Maroon},
  citecolor={Blue},
  urlcolor={Blue},
  pdfcreator={LaTeX via pandoc}}


\title{Social Network Analysis of the Canadian Disability Rights
Community}
\usepackage{etoolbox}
\makeatletter
\providecommand{\subtitle}[1]{% add subtitle to \maketitle
  \apptocmd{\@title}{\par {\large #1 \par}}{}{}
}
\makeatother
\subtitle{Mapping Collaboration Networks and Key Actors in Disability
Rights Advocacy}
\author{Dr.~Derrick L. Cogburn}
\date{2026-03-28}

\begin{document}
\maketitle

\renewcommand*\contentsname{Table of contents}
{
\hypersetup{linkcolor=}
\setcounter{tocdepth}{2}
\tableofcontents
}

\section{Executive Summary}\label{executive-summary}

This report presents a Social Network Analysis (SNA) of collaboration
patterns within the Canadian disability rights community, based on a
bilingual survey (English/French). The survey, called ``Enquête
d'analyse des Réseaux Sociaux / Social Network Analysis Survey'' enables
an analysis and mapping of who collaborates with whom, identifies the
most central actors and organizations, detects community clusters, and
assesses overall network structure.

\section{Data Cleaning and Pipeline}\label{data-cleaning-and-pipeline}

\subsection{Source Data}\label{source-data}

The raw data was collected via a bilingual (English/French) Microsoft
Forms survey distributed to members of the Canadian disability rights
community. The survey instrument contained 51 columns across a single
worksheet, with parallel English and French versions of each question. A
total of 111 responses were recorded. The raw data file was
de-identified prior to analysis.

\subsection{Survey Structure and Key
Fields}\label{survey-structure-and-key-fields}

The survey captured several categories of information relevant to social
network analysis. Respondent attributes included the respondent's
connection to the disability community (e.g., person with a disability,
family member/caregiver, advocate, researcher, health professional) and
the disability rights issues in which they are most actively involved
(free-text responses covering topics such as employment, children/youth,
housing, poverty, and accessibility). Collaboration data included the
names of up to three organizations or individuals with whom the
respondent has partnered in the past five years (free-text), the
frequency of collaboration with each named partner (regularly,
occasionally, or seldom), and whether the respondent's named partners
are linked to each other.

\subsection{Data Cleaning Steps}\label{data-cleaning-steps}

The cleaning pipeline proceeded through several stages. First, the
metadata row generated by Microsoft Forms was removed, and the language
preference field was used to route each respondent to the appropriate
English or French column for each question. Second, the bilingual fields
were consolidated, with the French-language response used when the
respondent selected French and the English-language response used
otherwise. This produced a unified respondent table with consistent
column names regardless of survey language.

Third, connection types were standardized into nine categories: Person
with Disability, Advocate/CSO, Family/Caregiver, Researcher, Health,
Education, Leisure, Government/Policy, and Other. Respondents whose
connection type text matched multiple categories were assigned all
applicable categories (semicolon-separated), with the first match
designated as the primary connection type.

Fourth, disability rights issue areas were categorized from the
free-text responses into 12 standardized topics: Children/Youth,
Employment, Accessibility/Inclusion, Poverty/Income, Rights/Policy,
Education, Housing, Health, Indigenous, Transportation, Technology, and
Other. As with connection types, respondents active in multiple areas
received all applicable categories.

Fifth, the free-text partner names were parsed into individual
organization or person names. The parser handled several formats
including comma-separated lists, semicolon-separated lists, numbered
lists, and narrative descriptions. Each respondent contributed up to
three partner names. Names were then normalized by converting to
lowercase, collapsing whitespace, and removing punctuation. A
prefix-matching algorithm was applied to merge near-duplicate entries
(e.g., variant spellings or abbreviations of the same organization).
This best-effort automated cleaning reduced the number of unique partner
names from 175 to 166.

Sixth, collaboration frequency values were harmonized across languages,
mapping both English terms (Regularly, Occasionally, Seldom) and their
French equivalents (Régulièrement, Occasionnellement, Rarement) to a
common three-level scale plus an Unknown category for non-standard
responses.

\subsection{Output Files}\label{output-files}

The pipeline produced two analysis-ready CSV files. The node table
(\texttt{respondents\_final.csv}) contains one row per respondent with
fields for respondent ID, language, raw and categorized connection type,
raw and categorized issue areas, and primary connection and issue
assignments. The edge table (\texttt{edges\_final.csv}) contains one row
per respondent-partner tie with fields for respondent ID, raw partner
name, normalized partner name, cleaned canonical partner name, raw
frequency, and harmonized frequency.

\section{Setup}\label{setup}

\subsection{Install Required Packages}\label{install-required-packages}

Run this chunk once to install all dependencies. After installation, you
can set \texttt{eval:\ false} or simply skip it on future renders.

\begin{Shaded}
\begin{Highlighting}[]
\FunctionTok{install.packages}\NormalTok{(}\FunctionTok{c}\NormalTok{(}
  \StringTok{"tidyverse"}\NormalTok{,      }\CommentTok{\# Data wrangling and ggplot2}
  \StringTok{"igraph"}\NormalTok{,         }\CommentTok{\# Core network analysis}
  \StringTok{"tidygraph"}\NormalTok{,      }\CommentTok{\# Tidy interface to igraph}
  \StringTok{"ggraph"}\NormalTok{,         }\CommentTok{\# Network visualisation (ggplot2{-}based)}
  \StringTok{"readxl"}\NormalTok{,         }\CommentTok{\# Read .xlsx files}
  \StringTok{"patchwork"}\NormalTok{,      }\CommentTok{\# Combine ggplot panels}
  \StringTok{"scales"}\NormalTok{,         }\CommentTok{\# Axis formatting}
  \StringTok{"RColorBrewer"}\NormalTok{,   }\CommentTok{\# Colour palettes}
  \StringTok{"knitr"}\NormalTok{,          }\CommentTok{\# Tables in Quarto}
  \StringTok{"kableExtra"}      \CommentTok{\# Styled tables}
\NormalTok{))}
\end{Highlighting}
\end{Shaded}

\subsection{Load Libraries}\label{load-libraries}

\begin{Shaded}
\begin{Highlighting}[]
\FunctionTok{library}\NormalTok{(tidyverse)}
\FunctionTok{library}\NormalTok{(igraph)}
\FunctionTok{library}\NormalTok{(tidygraph)}
\FunctionTok{library}\NormalTok{(ggraph)}
\FunctionTok{library}\NormalTok{(readxl)}
\FunctionTok{library}\NormalTok{(patchwork)}
\FunctionTok{library}\NormalTok{(scales)}
\FunctionTok{library}\NormalTok{(RColorBrewer)}
\FunctionTok{library}\NormalTok{(knitr)}
\FunctionTok{library}\NormalTok{(kableExtra)}

\FunctionTok{theme\_set}\NormalTok{(}\FunctionTok{theme\_minimal}\NormalTok{(}\AttributeTok{base\_size =} \DecValTok{12}\NormalTok{))}
\FunctionTok{set.seed}\NormalTok{(}\DecValTok{42}\NormalTok{)}
\end{Highlighting}
\end{Shaded}

\subsection{Import Data}\label{import-data}

The analysis uses two pre-cleaned CSV files produced by Dr.~Cogburn in a
data-preparation pipeline. These files, \texttt{respondents\_final.csv}
and \texttt{edges\_final.csv} are located in the data folder, which is
located in the same directory as this \texttt{.qmd} file. Paths to the
data should be updated to data/filename.

\begin{Shaded}
\begin{Highlighting}[]
\CommentTok{\# ── Node table: one row per survey respondent ──}
\NormalTok{respondents }\OtherTok{\textless{}{-}} \FunctionTok{read\_csv}\NormalTok{(}\StringTok{"data/respondents\_final.csv"}\NormalTok{, }\AttributeTok{show\_col\_types =} \ConstantTok{FALSE}\NormalTok{) }\SpecialCharTok{|\textgreater{}}
  \CommentTok{\# Drop the metadata row carried over from MS Forms}
\FunctionTok{filter}\NormalTok{(resp\_id }\SpecialCharTok{!=} \StringTok{"RMicrosoft Forms "}\NormalTok{)}

\CommentTok{\# ── Edge table: one row per respondent → partner tie ──}
\NormalTok{edges }\OtherTok{\textless{}{-}} \FunctionTok{read\_csv}\NormalTok{(}\StringTok{"data/edges\_final.csv"}\NormalTok{, }\AttributeTok{show\_col\_types =} \ConstantTok{FALSE}\NormalTok{) }\SpecialCharTok{|\textgreater{}}
  \FunctionTok{filter}\NormalTok{(}
    \SpecialCharTok{!}\FunctionTok{is.na}\NormalTok{(partner\_clean),}
\NormalTok{    partner\_clean }\SpecialCharTok{!=} \StringTok{"nan"}\NormalTok{,}
\NormalTok{    partner\_clean }\SpecialCharTok{!=} \StringTok{"None"}
\NormalTok{  )}
\end{Highlighting}
\end{Shaded}

\begin{Shaded}
\begin{Highlighting}[]
\FunctionTok{cat}\NormalTok{(}
  \StringTok{"Respondents:"}\NormalTok{, }\FunctionTok{nrow}\NormalTok{(respondents), }\StringTok{"}\SpecialCharTok{\textbackslash{}n}\StringTok{"}\NormalTok{,}
  \StringTok{"Respondents who named partners:"}\NormalTok{, }\FunctionTok{n\_distinct}\NormalTok{(edges}\SpecialCharTok{$}\NormalTok{respondent), }\StringTok{"}\SpecialCharTok{\textbackslash{}n}\StringTok{"}\NormalTok{,}
  \StringTok{"Unique organisations named:"}\NormalTok{, }\FunctionTok{n\_distinct}\NormalTok{(edges}\SpecialCharTok{$}\NormalTok{partner\_clean), }\StringTok{"}\SpecialCharTok{\textbackslash{}n}\StringTok{"}\NormalTok{,}
  \StringTok{"Total collaboration ties:"}\NormalTok{, }\FunctionTok{nrow}\NormalTok{(edges), }\StringTok{"}\SpecialCharTok{\textbackslash{}n}\StringTok{"}
\NormalTok{)}
\end{Highlighting}
\end{Shaded}

\begin{verbatim}
Respondents: 111 
 Respondents who named partners: 57 
 Unique organisations named: 166 
 Total collaboration ties: 182 
\end{verbatim}

\subsection{Colour Palettes}\label{colour-palettes}

\begin{Shaded}
\begin{Highlighting}[]
\CommentTok{\# Connection{-}type palette (respondent node colouring)}
\NormalTok{conn\_palette }\OtherTok{\textless{}{-}} \FunctionTok{c}\NormalTok{(}
  \StringTok{"Person with Disability"} \OtherTok{=} \StringTok{"\#E63946"}\NormalTok{,}
  \StringTok{"Advocate/CSO"}           \OtherTok{=} \StringTok{"\#457B9D"}\NormalTok{,}
  \StringTok{"Family/Caregiver"}       \OtherTok{=} \StringTok{"\#2A9D8F"}\NormalTok{,}
  \StringTok{"Researcher"}             \OtherTok{=} \StringTok{"\#6A4C93"}\NormalTok{,}
  \StringTok{"Health"}                 \OtherTok{=} \StringTok{"\#BC4749"}\NormalTok{,}
  \StringTok{"Education"}              \OtherTok{=} \StringTok{"\#F4A261"}\NormalTok{,}
  \StringTok{"Leisure"}                \OtherTok{=} \StringTok{"\#8AC926"}\NormalTok{,}
  \StringTok{"Government/Policy"}      \OtherTok{=} \StringTok{"\#264653"}\NormalTok{,}
  \StringTok{"Other"}                  \OtherTok{=} \StringTok{"\#AAAAAA"}\NormalTok{,}
  \StringTok{"Unknown"}                \OtherTok{=} \StringTok{"\#BBBBBB"}
\NormalTok{)}

\CommentTok{\# Issue{-}area palette (bipartite squares)}
\NormalTok{issue\_palette }\OtherTok{\textless{}{-}} \FunctionTok{c}\NormalTok{(}
  \StringTok{"Children/Youth"}          \OtherTok{=} \StringTok{"\#E63946"}\NormalTok{,}
  \StringTok{"Employment"}              \OtherTok{=} \StringTok{"\#457B9D"}\NormalTok{,}
  \StringTok{"Accessibility/Inclusion"} \OtherTok{=} \StringTok{"\#2A9D8F"}\NormalTok{,}
  \StringTok{"Poverty/Income"}          \OtherTok{=} \StringTok{"\#E9C46A"}\NormalTok{,}
  \StringTok{"Rights/Policy"}           \OtherTok{=} \StringTok{"\#264653"}\NormalTok{,}
  \StringTok{"Education"}               \OtherTok{=} \StringTok{"\#F4A261"}\NormalTok{,}
  \StringTok{"Housing"}                 \OtherTok{=} \StringTok{"\#606C38"}\NormalTok{,}
  \StringTok{"Health"}                  \OtherTok{=} \StringTok{"\#BC4749"}\NormalTok{,}
  \StringTok{"Indigenous"}              \OtherTok{=} \StringTok{"\#6A4C93"}\NormalTok{,}
  \StringTok{"Transportation"}          \OtherTok{=} \StringTok{"\#1982C4"}\NormalTok{,}
  \StringTok{"Technology"}              \OtherTok{=} \StringTok{"\#8AC926"}\NormalTok{,}
  \StringTok{"Other"}                   \OtherTok{=} \StringTok{"\#AAAAAA"}
\NormalTok{)}

\CommentTok{\# Collaboration{-}frequency palette}
\NormalTok{freq\_palette }\OtherTok{\textless{}{-}} \FunctionTok{c}\NormalTok{(}
  \StringTok{"Regularly"}    \OtherTok{=} \StringTok{"\#E63946"}\NormalTok{,}
  \StringTok{"Occasionally"} \OtherTok{=} \StringTok{"\#E9C46A"}\NormalTok{,}
  \StringTok{"Seldom"}       \OtherTok{=} \StringTok{"\#457B9D"}\NormalTok{,}
  \StringTok{"Unknown"}      \OtherTok{=} \StringTok{"\#BBBBBB"}
\NormalTok{)}
\end{Highlighting}
\end{Shaded}

\section{Two-Mode Bipartite Network: Respondents × Issue
Areas}\label{two-mode-bipartite-network-respondents-issue-areas}

This network maps respondents (circles) to the disability rights issue
areas they are most actively involved in (squares). Respondents are
coloured by their connection type to the disability community;
issue-area nodes are sized by the number of affiliated respondents.

\begin{Shaded}
\begin{Highlighting}[]
\CommentTok{\# Expand the semicolon{-}separated issue\_category into long form}
\NormalTok{bip\_edges }\OtherTok{\textless{}{-}}\NormalTok{ respondents }\SpecialCharTok{|\textgreater{}}
  \FunctionTok{filter}\NormalTok{(}
    \SpecialCharTok{!}\FunctionTok{is.na}\NormalTok{(issue\_category),}
\NormalTok{    issue\_category }\SpecialCharTok{!=} \StringTok{"Unknown"}
\NormalTok{  ) }\SpecialCharTok{|\textgreater{}}
  \FunctionTok{separate\_rows}\NormalTok{(issue\_category, }\AttributeTok{sep =} \StringTok{";"}\NormalTok{) }\SpecialCharTok{|\textgreater{}}
  \FunctionTok{mutate}\NormalTok{(}\AttributeTok{issue\_category =} \FunctionTok{str\_trim}\NormalTok{(issue\_category)) }\SpecialCharTok{|\textgreater{}}
  \FunctionTok{filter}\NormalTok{(issue\_category }\SpecialCharTok{!=} \StringTok{""}\NormalTok{, issue\_category }\SpecialCharTok{!=} \StringTok{"Unknown"}\NormalTok{) }\SpecialCharTok{|\textgreater{}}
  \FunctionTok{select}\NormalTok{(resp\_id, issue\_category)}

\CommentTok{\# Build igraph bipartite}
\NormalTok{bip\_graph }\OtherTok{\textless{}{-}} \FunctionTok{graph\_from\_data\_frame}\NormalTok{(bip\_edges, }\AttributeTok{directed =} \ConstantTok{FALSE}\NormalTok{)}

\CommentTok{\# Mark node type: TRUE = issue{-}area node, FALSE = respondent node}
\FunctionTok{V}\NormalTok{(bip\_graph)}\SpecialCharTok{$}\NormalTok{type }\OtherTok{\textless{}{-}} \FunctionTok{V}\NormalTok{(bip\_graph)}\SpecialCharTok{$}\NormalTok{name }\SpecialCharTok{\%in\%} \FunctionTok{unique}\NormalTok{(bip\_edges}\SpecialCharTok{$}\NormalTok{issue\_category)}

\CommentTok{\# Attach respondent attributes}
\NormalTok{resp\_lookup }\OtherTok{\textless{}{-}}\NormalTok{ respondents }\SpecialCharTok{|\textgreater{}} \FunctionTok{select}\NormalTok{(resp\_id, primary\_connection)}
\FunctionTok{V}\NormalTok{(bip\_graph)}\SpecialCharTok{$}\NormalTok{primary\_connection }\OtherTok{\textless{}{-}}\NormalTok{ resp\_lookup}\SpecialCharTok{$}\NormalTok{primary\_connection[}
  \FunctionTok{match}\NormalTok{(}\FunctionTok{V}\NormalTok{(bip\_graph)}\SpecialCharTok{$}\NormalTok{name, resp\_lookup}\SpecialCharTok{$}\NormalTok{resp\_id)}
\NormalTok{]}
\FunctionTok{V}\NormalTok{(bip\_graph)}\SpecialCharTok{$}\NormalTok{primary\_connection[}\FunctionTok{is.na}\NormalTok{(}\FunctionTok{V}\NormalTok{(bip\_graph)}\SpecialCharTok{$}\NormalTok{primary\_connection)] }\OtherTok{\textless{}{-}} \StringTok{"Issue Area"}
\end{Highlighting}
\end{Shaded}

\begin{Shaded}
\begin{Highlighting}[]
\FunctionTok{as\_tbl\_graph}\NormalTok{(bip\_graph) }\SpecialCharTok{|\textgreater{}}
  \FunctionTok{mutate}\NormalTok{(}
    \AttributeTok{node\_type   =} \FunctionTok{ifelse}\NormalTok{(type, }\StringTok{"Issue Area"}\NormalTok{, }\StringTok{"Respondent"}\NormalTok{),}
    \AttributeTok{node\_degree =} \FunctionTok{centrality\_degree}\NormalTok{(),}
    \AttributeTok{node\_colour =} \FunctionTok{ifelse}\NormalTok{(type, }\StringTok{"Issue Area"}\NormalTok{, primary\_connection),}
    \AttributeTok{node\_size   =} \FunctionTok{ifelse}\NormalTok{(type, node\_degree }\SpecialCharTok{*} \DecValTok{2} \SpecialCharTok{+} \DecValTok{4}\NormalTok{, node\_degree }\SpecialCharTok{+} \DecValTok{2}\NormalTok{)}
\NormalTok{  ) }\SpecialCharTok{|\textgreater{}}
  \FunctionTok{ggraph}\NormalTok{(}\AttributeTok{layout =} \StringTok{"fr"}\NormalTok{) }\SpecialCharTok{+}
  \FunctionTok{geom\_edge\_link}\NormalTok{(}\AttributeTok{alpha =} \FloatTok{0.07}\NormalTok{, }\AttributeTok{colour =} \StringTok{"\#666666"}\NormalTok{, }\AttributeTok{width =} \FloatTok{0.4}\NormalTok{) }\SpecialCharTok{+}
  \FunctionTok{geom\_node\_point}\NormalTok{(}
    \FunctionTok{aes}\NormalTok{(}\AttributeTok{colour =}\NormalTok{ node\_colour, }\AttributeTok{size =}\NormalTok{ node\_size, }\AttributeTok{shape =}\NormalTok{ node\_type),}
    \AttributeTok{alpha =} \FloatTok{0.75}
\NormalTok{  ) }\SpecialCharTok{+}
  \FunctionTok{geom\_node\_text}\NormalTok{(}
    \FunctionTok{aes}\NormalTok{(}\AttributeTok{label =} \FunctionTok{ifelse}\NormalTok{(type, name, }\ConstantTok{NA}\NormalTok{)),}
    \AttributeTok{size =} \DecValTok{3}\NormalTok{, }\AttributeTok{fontface =} \StringTok{"bold"}\NormalTok{, }\AttributeTok{repel =} \ConstantTok{TRUE}\NormalTok{, }\AttributeTok{max.overlaps =} \DecValTok{20}
\NormalTok{  ) }\SpecialCharTok{+}
  \FunctionTok{scale\_shape\_manual}\NormalTok{(}\AttributeTok{values =} \FunctionTok{c}\NormalTok{(}\StringTok{"Issue Area"} \OtherTok{=} \DecValTok{15}\NormalTok{, }\StringTok{"Respondent"} \OtherTok{=} \DecValTok{16}\NormalTok{)) }\SpecialCharTok{+}
  \FunctionTok{scale\_colour\_manual}\NormalTok{(}
    \AttributeTok{values =} \FunctionTok{c}\NormalTok{(conn\_palette, }\StringTok{"Issue Area"} \OtherTok{=} \StringTok{"\#333333"}\NormalTok{),}
    \AttributeTok{name   =} \StringTok{"Node Type"}
\NormalTok{  ) }\SpecialCharTok{+}
  \FunctionTok{scale\_size\_identity}\NormalTok{() }\SpecialCharTok{+}
  \FunctionTok{guides}\NormalTok{(}\AttributeTok{shape =} \StringTok{"none"}\NormalTok{, }\AttributeTok{size =} \StringTok{"none"}\NormalTok{) }\SpecialCharTok{+}
  \FunctionTok{labs}\NormalTok{(}\AttributeTok{title =} \StringTok{"Bipartite Network: Respondents and Disability Rights Issue Areas"}\NormalTok{) }\SpecialCharTok{+}
  \FunctionTok{theme\_void}\NormalTok{() }\SpecialCharTok{+}
  \FunctionTok{theme}\NormalTok{(}
    \AttributeTok{legend.position =} \StringTok{"bottom"}\NormalTok{,}
    \AttributeTok{plot.title =} \FunctionTok{element\_text}\NormalTok{(}\AttributeTok{face =} \StringTok{"bold"}\NormalTok{, }\AttributeTok{size =} \DecValTok{14}\NormalTok{)}
\NormalTok{  )}
\end{Highlighting}
\end{Shaded}

\begin{figure}[H]

{\centering \pandocbounded{\includegraphics[keepaspectratio]{sna_analysis-2_files/figure-pdf/bipartite-plot-1.pdf}}

}

\caption{Figure 1. Two-mode bipartite network. Circles = respondents
(coloured by connection type); squares = issue areas (sized by degree).}

\end{figure}%

\section{One-Mode Projected Network: Respondent
Collaboration}\label{one-mode-projected-network-respondent-collaboration}

Two respondents are linked in the projected network when they named the
same organisation as a collaboration partner. Edge weight reflects the
number of shared partners.

\begin{Shaded}
\begin{Highlighting}[]
\CommentTok{\# Build the respondent–organisation bipartite graph}
\NormalTok{resp\_org\_edges }\OtherTok{\textless{}{-}}\NormalTok{ edges }\SpecialCharTok{|\textgreater{}}
  \FunctionTok{select}\NormalTok{(respondent, partner\_clean) }\SpecialCharTok{|\textgreater{}}
  \FunctionTok{distinct}\NormalTok{()}

\NormalTok{bip\_ro }\OtherTok{\textless{}{-}} \FunctionTok{graph\_from\_data\_frame}\NormalTok{(resp\_org\_edges, }\AttributeTok{directed =} \ConstantTok{FALSE}\NormalTok{)}
\FunctionTok{V}\NormalTok{(bip\_ro)}\SpecialCharTok{$}\NormalTok{type }\OtherTok{\textless{}{-}} \SpecialCharTok{!}\NormalTok{(}\FunctionTok{V}\NormalTok{(bip\_ro)}\SpecialCharTok{$}\NormalTok{name }\SpecialCharTok{\%in\%} \FunctionTok{unique}\NormalTok{(resp\_org\_edges}\SpecialCharTok{$}\NormalTok{respondent))}

\CommentTok{\# Project onto respondent mode}
\NormalTok{proj }\OtherTok{\textless{}{-}} \FunctionTok{bipartite\_projection}\NormalTok{(bip\_ro, }\AttributeTok{which =} \StringTok{"false"}\NormalTok{)}

\CommentTok{\# Attach respondent attributes}
\FunctionTok{V}\NormalTok{(proj)}\SpecialCharTok{$}\NormalTok{primary\_connection }\OtherTok{\textless{}{-}}\NormalTok{ resp\_lookup}\SpecialCharTok{$}\NormalTok{primary\_connection[}
  \FunctionTok{match}\NormalTok{(}\FunctionTok{V}\NormalTok{(proj)}\SpecialCharTok{$}\NormalTok{name, resp\_lookup}\SpecialCharTok{$}\NormalTok{resp\_id)}
\NormalTok{]}
\FunctionTok{V}\NormalTok{(proj)}\SpecialCharTok{$}\NormalTok{primary\_connection[}\FunctionTok{is.na}\NormalTok{(}\FunctionTok{V}\NormalTok{(proj)}\SpecialCharTok{$}\NormalTok{primary\_connection)] }\OtherTok{\textless{}{-}} \StringTok{"Unknown"}

\NormalTok{issue\_lookup }\OtherTok{\textless{}{-}}\NormalTok{ respondents }\SpecialCharTok{|\textgreater{}} \FunctionTok{select}\NormalTok{(resp\_id, primary\_issue)}
\FunctionTok{V}\NormalTok{(proj)}\SpecialCharTok{$}\NormalTok{primary\_issue }\OtherTok{\textless{}{-}}\NormalTok{ issue\_lookup}\SpecialCharTok{$}\NormalTok{primary\_issue[}
  \FunctionTok{match}\NormalTok{(}\FunctionTok{V}\NormalTok{(proj)}\SpecialCharTok{$}\NormalTok{name, issue\_lookup}\SpecialCharTok{$}\NormalTok{resp\_id)}
\NormalTok{]}
\end{Highlighting}
\end{Shaded}

\begin{Shaded}
\begin{Highlighting}[]
\FunctionTok{as\_tbl\_graph}\NormalTok{(proj) }\SpecialCharTok{|\textgreater{}}
  \FunctionTok{mutate}\NormalTok{(}\AttributeTok{deg =} \FunctionTok{centrality\_degree}\NormalTok{()) }\SpecialCharTok{|\textgreater{}}
  \FunctionTok{ggraph}\NormalTok{(}\AttributeTok{layout =} \StringTok{"fr"}\NormalTok{) }\SpecialCharTok{+}
  \FunctionTok{geom\_edge\_link}\NormalTok{(}
    \FunctionTok{aes}\NormalTok{(}\AttributeTok{width =}\NormalTok{ weight),}
    \AttributeTok{alpha =} \FloatTok{0.25}\NormalTok{, }\AttributeTok{colour =} \StringTok{"\#999999"}
\NormalTok{  ) }\SpecialCharTok{+}
  \FunctionTok{geom\_node\_point}\NormalTok{(}
    \FunctionTok{aes}\NormalTok{(}\AttributeTok{colour =}\NormalTok{ primary\_connection, }\AttributeTok{size =}\NormalTok{ deg),}
    \AttributeTok{alpha =} \FloatTok{0.8}
\NormalTok{  ) }\SpecialCharTok{+}
  \FunctionTok{geom\_node\_text}\NormalTok{(}
    \FunctionTok{aes}\NormalTok{(}\AttributeTok{label =} \FunctionTok{ifelse}\NormalTok{(deg }\SpecialCharTok{\textgreater{}=} \DecValTok{2}\NormalTok{, name, }\ConstantTok{NA}\NormalTok{)),}
    \AttributeTok{size =} \FloatTok{2.5}\NormalTok{, }\AttributeTok{repel =} \ConstantTok{TRUE}\NormalTok{, }\AttributeTok{max.overlaps =} \DecValTok{15}
\NormalTok{  ) }\SpecialCharTok{+}
  \FunctionTok{scale\_edge\_width}\NormalTok{(}\AttributeTok{range =} \FunctionTok{c}\NormalTok{(}\FloatTok{0.5}\NormalTok{, }\DecValTok{3}\NormalTok{), }\AttributeTok{name =} \StringTok{"Shared Partners"}\NormalTok{) }\SpecialCharTok{+}
  \FunctionTok{scale\_colour\_manual}\NormalTok{(}\AttributeTok{values =}\NormalTok{ conn\_palette, }\AttributeTok{name =} \StringTok{"Connection Type"}\NormalTok{) }\SpecialCharTok{+}
  \FunctionTok{scale\_size\_continuous}\NormalTok{(}\AttributeTok{range =} \FunctionTok{c}\NormalTok{(}\DecValTok{3}\NormalTok{, }\DecValTok{14}\NormalTok{), }\AttributeTok{name =} \StringTok{"Degree"}\NormalTok{) }\SpecialCharTok{+}
  \FunctionTok{labs}\NormalTok{(}\AttributeTok{title =} \StringTok{"Projected Network: Respondents Linked by Shared Collaboration Partners"}\NormalTok{) }\SpecialCharTok{+}
  \FunctionTok{theme\_void}\NormalTok{() }\SpecialCharTok{+}
  \FunctionTok{theme}\NormalTok{(}
    \AttributeTok{legend.position =} \StringTok{"bottom"}\NormalTok{,}
    \AttributeTok{plot.title =} \FunctionTok{element\_text}\NormalTok{(}\AttributeTok{face =} \StringTok{"bold"}\NormalTok{, }\AttributeTok{size =} \DecValTok{14}\NormalTok{)}
\NormalTok{  )}
\end{Highlighting}
\end{Shaded}

\begin{figure}[H]

{\centering \pandocbounded{\includegraphics[keepaspectratio]{sna_analysis-2_files/figure-pdf/projected-plot-1.pdf}}

}

\caption{Figure 2. One-mode projected network. Node size = degree;
colour = connection type; edge width = number of shared partners.}

\end{figure}%

\subsection{Network Structure Metrics}\label{network-structure-metrics}

\begin{Shaded}
\begin{Highlighting}[]
\NormalTok{metrics\_tbl }\OtherTok{\textless{}{-}} \FunctionTok{tibble}\NormalTok{(}
  \AttributeTok{Metric =} \FunctionTok{c}\NormalTok{(}
    \StringTok{"Nodes (respondents with ties)"}\NormalTok{,}
    \StringTok{"Edges (shared{-}partner links)"}\NormalTok{,}
    \StringTok{"Density"}\NormalTok{,}
    \StringTok{"Connected components"}\NormalTok{,}
    \StringTok{"Transitivity (global clustering)"}\NormalTok{,}
    \StringTok{"Average clustering coefficient"}\NormalTok{,}
    \StringTok{"Largest component size"}
\NormalTok{  ),}
  \AttributeTok{Value =} \FunctionTok{c}\NormalTok{(}
    \FunctionTok{vcount}\NormalTok{(proj),}
    \FunctionTok{ecount}\NormalTok{(proj),}
    \FunctionTok{round}\NormalTok{(}\FunctionTok{graph.density}\NormalTok{(proj), }\DecValTok{4}\NormalTok{),}
    \FunctionTok{components}\NormalTok{(proj)}\SpecialCharTok{$}\NormalTok{no,}
    \FunctionTok{round}\NormalTok{(}\FunctionTok{transitivity}\NormalTok{(proj, }\AttributeTok{type =} \StringTok{"global"}\NormalTok{), }\DecValTok{4}\NormalTok{),}
    \FunctionTok{round}\NormalTok{(}\FunctionTok{transitivity}\NormalTok{(proj, }\AttributeTok{type =} \StringTok{"average"}\NormalTok{), }\DecValTok{4}\NormalTok{),}
    \FunctionTok{max}\NormalTok{(}\FunctionTok{components}\NormalTok{(proj)}\SpecialCharTok{$}\NormalTok{csize)}
\NormalTok{  )}
\NormalTok{)}

\NormalTok{metrics\_tbl }\SpecialCharTok{|\textgreater{}}
  \FunctionTok{kable}\NormalTok{(}\AttributeTok{align =} \FunctionTok{c}\NormalTok{(}\StringTok{"l"}\NormalTok{, }\StringTok{"r"}\NormalTok{), }\AttributeTok{booktabs =} \ConstantTok{TRUE}\NormalTok{) }\SpecialCharTok{|\textgreater{}}
  \FunctionTok{kable\_styling}\NormalTok{(}
    \AttributeTok{bootstrap\_options =} \FunctionTok{c}\NormalTok{(}\StringTok{"striped"}\NormalTok{, }\StringTok{"hover"}\NormalTok{),}
    \AttributeTok{latex\_options =} \FunctionTok{c}\NormalTok{(}\StringTok{"hold\_position"}\NormalTok{),}
    \AttributeTok{full\_width =} \ConstantTok{FALSE}
\NormalTok{  )}
\end{Highlighting}
\end{Shaded}

\begin{longtable*}[t]{lr}
\toprule
Metric & Value\\
\midrule
Nodes (respondents with ties) & 57.0000\\
Edges (shared-partner links) & 17.0000\\
Density & 0.0107\\
Connected components & 46.0000\\
Transitivity (global clustering) & 0.4773\\
\addlinespace
Average clustering coefficient & 0.8116\\
Largest component size & 10.0000\\
\bottomrule
\end{longtable*}

\section{Centrality Analysis}\label{centrality-analysis}

Centrality measures identify the most influential actors. \textbf{Degree
centrality} counts direct connections; \textbf{betweenness centrality}
identifies actors who bridge otherwise disconnected parts of the
network.

\subsection{Organisation Degree
Centrality}\label{organisation-degree-centrality}

\begin{Shaded}
\begin{Highlighting}[]
\NormalTok{org\_degree }\OtherTok{\textless{}{-}}\NormalTok{ edges }\SpecialCharTok{|\textgreater{}}
  \FunctionTok{count}\NormalTok{(partner\_clean, }\AttributeTok{sort =} \ConstantTok{TRUE}\NormalTok{, }\AttributeTok{name =} \StringTok{"degree"}\NormalTok{) }\SpecialCharTok{|\textgreater{}}
  \FunctionTok{head}\NormalTok{(}\DecValTok{15}\NormalTok{)}

\FunctionTok{ggplot}\NormalTok{(org\_degree, }\FunctionTok{aes}\NormalTok{(}\AttributeTok{x =} \FunctionTok{reorder}\NormalTok{(partner\_clean, degree), }\AttributeTok{y =}\NormalTok{ degree)) }\SpecialCharTok{+}
  \FunctionTok{geom\_col}\NormalTok{(}\AttributeTok{fill =} \StringTok{"\#457B9D"}\NormalTok{, }\AttributeTok{alpha =} \FloatTok{0.85}\NormalTok{, }\AttributeTok{width =} \FloatTok{0.7}\NormalTok{) }\SpecialCharTok{+}
  \FunctionTok{geom\_text}\NormalTok{(}\FunctionTok{aes}\NormalTok{(}\AttributeTok{label =}\NormalTok{ degree), }\AttributeTok{hjust =} \SpecialCharTok{{-}}\FloatTok{0.3}\NormalTok{, }\AttributeTok{size =} \FloatTok{3.5}\NormalTok{) }\SpecialCharTok{+}
  \FunctionTok{coord\_flip}\NormalTok{() }\SpecialCharTok{+}
  \FunctionTok{labs}\NormalTok{(}
    \AttributeTok{x =} \ConstantTok{NULL}\NormalTok{,}
    \AttributeTok{y =} \StringTok{"Degree (respondents naming this organisation)"}\NormalTok{,}
    \AttributeTok{title =} \StringTok{"Organisation Degree Centrality"}
\NormalTok{  ) }\SpecialCharTok{+}
  \FunctionTok{theme}\NormalTok{(}\AttributeTok{plot.title =} \FunctionTok{element\_text}\NormalTok{(}\AttributeTok{face =} \StringTok{"bold"}\NormalTok{))}
\end{Highlighting}
\end{Shaded}

\begin{figure}[H]

{\centering \pandocbounded{\includegraphics[keepaspectratio]{sna_analysis-2_files/figure-pdf/org-degree-1.pdf}}

}

\caption{Figure 3a. The 15 most frequently named collaboration
partners.}

\end{figure}%

\subsection{Respondent Centrality (Projected
Network)}\label{respondent-centrality-projected-network}

\begin{Shaded}
\begin{Highlighting}[]
\CommentTok{\# Compute centrality measures}
\NormalTok{cent\_df }\OtherTok{\textless{}{-}} \FunctionTok{tibble}\NormalTok{(}
  \AttributeTok{name              =} \FunctionTok{V}\NormalTok{(proj)}\SpecialCharTok{$}\NormalTok{name,}
  \AttributeTok{primary\_connection =} \FunctionTok{V}\NormalTok{(proj)}\SpecialCharTok{$}\NormalTok{primary\_connection,}
  \AttributeTok{degree            =} \FunctionTok{degree}\NormalTok{(proj),}
  \AttributeTok{betweenness       =} \FunctionTok{betweenness}\NormalTok{(proj, }\AttributeTok{weights =} \FunctionTok{E}\NormalTok{(proj)}\SpecialCharTok{$}\NormalTok{weight)}
\NormalTok{) }\SpecialCharTok{|\textgreater{}}
  \FunctionTok{arrange}\NormalTok{(}\FunctionTok{desc}\NormalTok{(degree))}

\CommentTok{\# Degree panel}
\NormalTok{p\_deg }\OtherTok{\textless{}{-}}\NormalTok{ cent\_df }\SpecialCharTok{|\textgreater{}}
  \FunctionTok{head}\NormalTok{(}\DecValTok{12}\NormalTok{) }\SpecialCharTok{|\textgreater{}}
  \FunctionTok{mutate}\NormalTok{(}\AttributeTok{label =} \FunctionTok{paste0}\NormalTok{(name, }\StringTok{" ("}\NormalTok{, }\FunctionTok{str\_sub}\NormalTok{(primary\_connection, }\DecValTok{1}\NormalTok{, }\DecValTok{12}\NormalTok{), }\StringTok{")"}\NormalTok{)) }\SpecialCharTok{|\textgreater{}}
  \FunctionTok{ggplot}\NormalTok{(}\FunctionTok{aes}\NormalTok{(}\AttributeTok{x =} \FunctionTok{reorder}\NormalTok{(label, degree), }\AttributeTok{y =}\NormalTok{ degree, }\AttributeTok{fill =}\NormalTok{ primary\_connection)) }\SpecialCharTok{+}
  \FunctionTok{geom\_col}\NormalTok{(}\AttributeTok{alpha =} \FloatTok{0.85}\NormalTok{, }\AttributeTok{width =} \FloatTok{0.7}\NormalTok{) }\SpecialCharTok{+}
  \FunctionTok{coord\_flip}\NormalTok{() }\SpecialCharTok{+}
  \FunctionTok{scale\_fill\_manual}\NormalTok{(}\AttributeTok{values =}\NormalTok{ conn\_palette) }\SpecialCharTok{+}
  \FunctionTok{labs}\NormalTok{(}\AttributeTok{x =} \ConstantTok{NULL}\NormalTok{, }\AttributeTok{y =} \StringTok{"Degree"}\NormalTok{, }\AttributeTok{title =} \StringTok{"Respondent Degree Centrality"}\NormalTok{) }\SpecialCharTok{+}
  \FunctionTok{theme}\NormalTok{(}
    \AttributeTok{legend.position =} \StringTok{"none"}\NormalTok{,}
    \AttributeTok{plot.title =} \FunctionTok{element\_text}\NormalTok{(}\AttributeTok{face =} \StringTok{"bold"}\NormalTok{)}
\NormalTok{  )}

\CommentTok{\# Betweenness panel}
\NormalTok{p\_betw }\OtherTok{\textless{}{-}}\NormalTok{ cent\_df }\SpecialCharTok{|\textgreater{}}
  \FunctionTok{arrange}\NormalTok{(}\FunctionTok{desc}\NormalTok{(betweenness)) }\SpecialCharTok{|\textgreater{}}
  \FunctionTok{head}\NormalTok{(}\DecValTok{12}\NormalTok{) }\SpecialCharTok{|\textgreater{}}
  \FunctionTok{mutate}\NormalTok{(}\AttributeTok{label =} \FunctionTok{paste0}\NormalTok{(name, }\StringTok{" ("}\NormalTok{, }\FunctionTok{str\_sub}\NormalTok{(primary\_connection, }\DecValTok{1}\NormalTok{, }\DecValTok{12}\NormalTok{), }\StringTok{")"}\NormalTok{)) }\SpecialCharTok{|\textgreater{}}
  \FunctionTok{ggplot}\NormalTok{(}\FunctionTok{aes}\NormalTok{(}\AttributeTok{x =} \FunctionTok{reorder}\NormalTok{(label, betweenness), }\AttributeTok{y =}\NormalTok{ betweenness, }\AttributeTok{fill =}\NormalTok{ primary\_connection)) }\SpecialCharTok{+}
  \FunctionTok{geom\_col}\NormalTok{(}\AttributeTok{alpha =} \FloatTok{0.85}\NormalTok{, }\AttributeTok{width =} \FloatTok{0.7}\NormalTok{) }\SpecialCharTok{+}
  \FunctionTok{coord\_flip}\NormalTok{() }\SpecialCharTok{+}
  \FunctionTok{scale\_fill\_manual}\NormalTok{(}\AttributeTok{values =}\NormalTok{ conn\_palette) }\SpecialCharTok{+}
  \FunctionTok{labs}\NormalTok{(}\AttributeTok{x =} \ConstantTok{NULL}\NormalTok{, }\AttributeTok{y =} \StringTok{"Betweenness"}\NormalTok{, }\AttributeTok{title =} \StringTok{"Respondent Betweenness Centrality"}\NormalTok{) }\SpecialCharTok{+}
  \FunctionTok{theme}\NormalTok{(}
    \AttributeTok{legend.position =} \StringTok{"none"}\NormalTok{,}
    \AttributeTok{plot.title =} \FunctionTok{element\_text}\NormalTok{(}\AttributeTok{face =} \StringTok{"bold"}\NormalTok{)}
\NormalTok{  )}

\NormalTok{p\_deg }\SpecialCharTok{+}\NormalTok{ p\_betw }\SpecialCharTok{+}
  \FunctionTok{plot\_annotation}\NormalTok{(}
    \AttributeTok{title =} \StringTok{"Centrality Analysis: Key Actors in the Projected Network"}\NormalTok{,}
    \AttributeTok{theme =} \FunctionTok{theme}\NormalTok{(}\AttributeTok{plot.title =} \FunctionTok{element\_text}\NormalTok{(}\AttributeTok{face =} \StringTok{"bold"}\NormalTok{, }\AttributeTok{size =} \DecValTok{14}\NormalTok{))}
\NormalTok{  )}
\end{Highlighting}
\end{Shaded}

\begin{figure}[H]

{\centering \includegraphics[width=1\linewidth,height=\textheight,keepaspectratio]{sna_analysis-2_files/figure-pdf/respondent-centrality-1.pdf}

}

\caption{Figure 3b. Respondent degree and betweenness centrality in the
projected network.}

\end{figure}%

\section{Community Detection}\label{community-detection}

The Louvain algorithm partitions the network into communities that
maximise modularity --- groups of respondents who are more densely
connected to each other than to the rest of the network.

\begin{Shaded}
\begin{Highlighting}[]
\NormalTok{comms }\OtherTok{\textless{}{-}} \FunctionTok{cluster\_louvain}\NormalTok{(proj, }\AttributeTok{weights =} \FunctionTok{E}\NormalTok{(proj)}\SpecialCharTok{$}\NormalTok{weight)}

\CommentTok{\# Summary stats}
\NormalTok{n\_comms    }\OtherTok{\textless{}{-}} \FunctionTok{length}\NormalTok{(}\FunctionTok{unique}\NormalTok{(}\FunctionTok{membership}\NormalTok{(comms)))}
\NormalTok{mod\_score  }\OtherTok{\textless{}{-}} \FunctionTok{round}\NormalTok{(}\FunctionTok{modularity}\NormalTok{(comms), }\DecValTok{4}\NormalTok{)}
\NormalTok{comm\_sizes }\OtherTok{\textless{}{-}} \FunctionTok{sort}\NormalTok{(}\FunctionTok{table}\NormalTok{(}\FunctionTok{membership}\NormalTok{(comms)), }\AttributeTok{decreasing =} \ConstantTok{TRUE}\NormalTok{)}

\FunctionTok{cat}\NormalTok{(}
  \StringTok{"Communities detected:"}\NormalTok{, n\_comms, }\StringTok{"}\SpecialCharTok{\textbackslash{}n}\StringTok{"}\NormalTok{,}
  \StringTok{"Modularity:"}\NormalTok{, mod\_score, }\StringTok{"}\SpecialCharTok{\textbackslash{}n}\StringTok{"}\NormalTok{,}
  \StringTok{"Sizes:"}\NormalTok{, }\FunctionTok{paste}\NormalTok{(comm\_sizes, }\AttributeTok{collapse =} \StringTok{", "}\NormalTok{), }\StringTok{"}\SpecialCharTok{\textbackslash{}n}\StringTok{"}
\NormalTok{)}
\end{Highlighting}
\end{Shaded}

\begin{verbatim}
Communities detected: 47 
 Modularity: 0.3387 
 Sizes: 7, 3, 2, 2, 1, 1, 1, 1, 1, 1, 1, 1, 1, 1, 1, 1, 1, 1, 1, 1, 1, 1, 1, 1, 1, 1, 1, 1, 1, 1, 1, 1, 1, 1, 1, 1, 1, 1, 1, 1, 1, 1, 1, 1, 1, 1, 1 
\end{verbatim}

\begin{Shaded}
\begin{Highlighting}[]
\FunctionTok{as\_tbl\_graph}\NormalTok{(proj) }\SpecialCharTok{|\textgreater{}}
  \FunctionTok{mutate}\NormalTok{(}
    \AttributeTok{community =} \FunctionTok{as.factor}\NormalTok{(}\FunctionTok{membership}\NormalTok{(comms)),}
    \AttributeTok{deg       =} \FunctionTok{centrality\_degree}\NormalTok{()}
\NormalTok{  ) }\SpecialCharTok{|\textgreater{}}
  \FunctionTok{ggraph}\NormalTok{(}\AttributeTok{layout =} \StringTok{"fr"}\NormalTok{) }\SpecialCharTok{+}
  \FunctionTok{geom\_edge\_link}\NormalTok{(}\AttributeTok{alpha =} \FloatTok{0.15}\NormalTok{, }\AttributeTok{colour =} \StringTok{"\#BBBBBB"}\NormalTok{, }\AttributeTok{width =} \FloatTok{0.5}\NormalTok{) }\SpecialCharTok{+}
  \FunctionTok{geom\_node\_point}\NormalTok{(}
    \FunctionTok{aes}\NormalTok{(}\AttributeTok{colour =}\NormalTok{ community, }\AttributeTok{size =}\NormalTok{ deg),}
    \AttributeTok{alpha =} \FloatTok{0.85}
\NormalTok{  ) }\SpecialCharTok{+}
  \FunctionTok{geom\_node\_text}\NormalTok{(}
    \FunctionTok{aes}\NormalTok{(}\AttributeTok{label =} \FunctionTok{ifelse}\NormalTok{(deg }\SpecialCharTok{\textgreater{}=} \DecValTok{2}\NormalTok{, name, }\ConstantTok{NA}\NormalTok{)),}
    \AttributeTok{size =} \FloatTok{2.5}\NormalTok{, }\AttributeTok{repel =} \ConstantTok{TRUE}\NormalTok{, }\AttributeTok{max.overlaps =} \DecValTok{15}
\NormalTok{  ) }\SpecialCharTok{+}
  \FunctionTok{scale\_size\_continuous}\NormalTok{(}\AttributeTok{range =} \FunctionTok{c}\NormalTok{(}\DecValTok{3}\NormalTok{, }\DecValTok{14}\NormalTok{), }\AttributeTok{name =} \StringTok{"Degree"}\NormalTok{) }\SpecialCharTok{+}
  \FunctionTok{labs}\NormalTok{(}
    \AttributeTok{title   =} \FunctionTok{paste0}\NormalTok{(}\StringTok{"Community Detection (Louvain): "}\NormalTok{, n\_comms, }\StringTok{" Communities"}\NormalTok{),}
    \AttributeTok{subtitle =} \FunctionTok{paste0}\NormalTok{(}\StringTok{"Modularity = "}\NormalTok{, mod\_score),}
    \AttributeTok{colour  =} \StringTok{"Community"}
\NormalTok{  ) }\SpecialCharTok{+}
  \FunctionTok{theme\_void}\NormalTok{() }\SpecialCharTok{+}
  \FunctionTok{theme}\NormalTok{(}
    \AttributeTok{legend.position =} \StringTok{"bottom"}\NormalTok{,}
    \AttributeTok{plot.title    =} \FunctionTok{element\_text}\NormalTok{(}\AttributeTok{face =} \StringTok{"bold"}\NormalTok{, }\AttributeTok{size =} \DecValTok{14}\NormalTok{),}
    \AttributeTok{plot.subtitle =} \FunctionTok{element\_text}\NormalTok{(}\AttributeTok{size =} \DecValTok{11}\NormalTok{, }\AttributeTok{colour =} \StringTok{"grey40"}\NormalTok{)}
\NormalTok{  )}
\end{Highlighting}
\end{Shaded}

\begin{figure}[H]

{\centering \pandocbounded{\includegraphics[keepaspectratio]{sna_analysis-2_files/figure-pdf/community-plot-1.pdf}}

}

\caption{Figure 4. Louvain communities. Colour = community membership;
size = degree.}

\end{figure}%

\subsection{Community Composition}\label{community-composition}

\begin{Shaded}
\begin{Highlighting}[]
\CommentTok{\# Build composition summary for multi{-}member communities}
\NormalTok{comm\_membership }\OtherTok{\textless{}{-}} \FunctionTok{tibble}\NormalTok{(}
  \AttributeTok{name              =} \FunctionTok{V}\NormalTok{(proj)}\SpecialCharTok{$}\NormalTok{name,}
  \AttributeTok{primary\_connection =} \FunctionTok{V}\NormalTok{(proj)}\SpecialCharTok{$}\NormalTok{primary\_connection,}
  \AttributeTok{primary\_issue     =} \FunctionTok{V}\NormalTok{(proj)}\SpecialCharTok{$}\NormalTok{primary\_issue,}
  \AttributeTok{community         =} \FunctionTok{membership}\NormalTok{(comms)}
\NormalTok{)}

\NormalTok{comm\_summary }\OtherTok{\textless{}{-}}\NormalTok{ comm\_membership }\SpecialCharTok{|\textgreater{}}
  \FunctionTok{group\_by}\NormalTok{(community) }\SpecialCharTok{|\textgreater{}}
  \FunctionTok{summarise}\NormalTok{(}
    \AttributeTok{Size =} \FunctionTok{n}\NormalTok{(),}
    \StringTok{\textasciigrave{}}\AttributeTok{Connection Types}\StringTok{\textasciigrave{}} \OtherTok{=} \FunctionTok{paste}\NormalTok{(}
      \FunctionTok{names}\NormalTok{(}\FunctionTok{sort}\NormalTok{(}\FunctionTok{table}\NormalTok{(primary\_connection), }\AttributeTok{decreasing =} \ConstantTok{TRUE}\NormalTok{)[}\DecValTok{1}\SpecialCharTok{:}\FunctionTok{min}\NormalTok{(}\DecValTok{3}\NormalTok{, }\FunctionTok{n}\NormalTok{())]),}
      \AttributeTok{collapse =} \StringTok{", "}
\NormalTok{    ),}
    \StringTok{\textasciigrave{}}\AttributeTok{Issue Areas}\StringTok{\textasciigrave{}} \OtherTok{=} \FunctionTok{paste}\NormalTok{(}
      \FunctionTok{names}\NormalTok{(}\FunctionTok{sort}\NormalTok{(}\FunctionTok{table}\NormalTok{(primary\_issue), }\AttributeTok{decreasing =} \ConstantTok{TRUE}\NormalTok{)[}\DecValTok{1}\SpecialCharTok{:}\FunctionTok{min}\NormalTok{(}\DecValTok{3}\NormalTok{, }\FunctionTok{n}\NormalTok{())]),}
      \AttributeTok{collapse =} \StringTok{", "}
\NormalTok{    ),}
    \AttributeTok{.groups =} \StringTok{"drop"}
\NormalTok{  ) }\SpecialCharTok{|\textgreater{}}
  \FunctionTok{filter}\NormalTok{(Size }\SpecialCharTok{\textgreater{}=} \DecValTok{2}\NormalTok{) }\SpecialCharTok{|\textgreater{}}
  \FunctionTok{arrange}\NormalTok{(}\FunctionTok{desc}\NormalTok{(Size)) }\SpecialCharTok{|\textgreater{}}
  \FunctionTok{mutate}\NormalTok{(}\AttributeTok{Community =} \FunctionTok{paste0}\NormalTok{(}\StringTok{"Community "}\NormalTok{, }\FunctionTok{row\_number}\NormalTok{())) }\SpecialCharTok{|\textgreater{}}
  \FunctionTok{select}\NormalTok{(Community, Size, }\StringTok{\textasciigrave{}}\AttributeTok{Connection Types}\StringTok{\textasciigrave{}}\NormalTok{, }\StringTok{\textasciigrave{}}\AttributeTok{Issue Areas}\StringTok{\textasciigrave{}}\NormalTok{)}

\NormalTok{comm\_summary }\SpecialCharTok{|\textgreater{}}
  \FunctionTok{kable}\NormalTok{(}\AttributeTok{booktabs =} \ConstantTok{TRUE}\NormalTok{) }\SpecialCharTok{|\textgreater{}}
  \FunctionTok{kable\_styling}\NormalTok{(}
    \AttributeTok{bootstrap\_options =} \FunctionTok{c}\NormalTok{(}\StringTok{"striped"}\NormalTok{, }\StringTok{"hover"}\NormalTok{),}
    \AttributeTok{latex\_options =} \FunctionTok{c}\NormalTok{(}\StringTok{"hold\_position"}\NormalTok{, }\StringTok{"scale\_down"}\NormalTok{),}
    \AttributeTok{full\_width =} \ConstantTok{FALSE}
\NormalTok{  )}
\end{Highlighting}
\end{Shaded}

\begin{longtable*}[t]{lrll}
\toprule
Community & Size & Connection Types & Issue Areas\\
\midrule
Community 1 & 7 & Person with Disability, Health, NA & Children/Youth, Employment, Housing\\
Community 2 & 3 & Person with Disability, NA, NA & Employment, Housing, Other\\
Community 3 & 2 & Person with Disability, NA & Employment, Poverty/Income\\
Community 4 & 2 & Person with Disability, NA & Children/Youth, NA\\
\bottomrule
\end{longtable*}

\section{Degree Distributions and Collaboration
Patterns}\label{degree-distributions-and-collaboration-patterns}

\begin{Shaded}
\begin{Highlighting}[]
\CommentTok{\# Organisation degree distribution}
\NormalTok{p1 }\OtherTok{\textless{}{-}}\NormalTok{ edges }\SpecialCharTok{|\textgreater{}}
  \FunctionTok{count}\NormalTok{(partner\_clean, }\AttributeTok{name =} \StringTok{"degree"}\NormalTok{) }\SpecialCharTok{|\textgreater{}}
  \FunctionTok{ggplot}\NormalTok{(}\FunctionTok{aes}\NormalTok{(}\AttributeTok{x =}\NormalTok{ degree)) }\SpecialCharTok{+}
  \FunctionTok{geom\_histogram}\NormalTok{(}
    \AttributeTok{binwidth =} \DecValTok{1}\NormalTok{, }\AttributeTok{fill =} \StringTok{"\#457B9D"}\NormalTok{, }\AttributeTok{alpha =} \FloatTok{0.8}\NormalTok{,}
    \AttributeTok{colour =} \StringTok{"white"}\NormalTok{, }\AttributeTok{boundary =} \FloatTok{0.5}
\NormalTok{  ) }\SpecialCharTok{+}
  \FunctionTok{geom\_vline}\NormalTok{(}
    \FunctionTok{aes}\NormalTok{(}\AttributeTok{xintercept =} \FunctionTok{mean}\NormalTok{(degree)),}
    \AttributeTok{linetype =} \StringTok{"dashed"}\NormalTok{, }\AttributeTok{colour =} \StringTok{"\#E63946"}
\NormalTok{  ) }\SpecialCharTok{+}
  \FunctionTok{labs}\NormalTok{(}
    \AttributeTok{x =} \StringTok{"Degree"}\NormalTok{, }\AttributeTok{y =} \StringTok{"Count"}\NormalTok{,}
    \AttributeTok{title =} \StringTok{"Organisation Degree Distribution"}
\NormalTok{  ) }\SpecialCharTok{+}
  \FunctionTok{theme}\NormalTok{(}\AttributeTok{plot.title =} \FunctionTok{element\_text}\NormalTok{(}\AttributeTok{face =} \StringTok{"bold"}\NormalTok{, }\AttributeTok{size =} \DecValTok{10}\NormalTok{))}

\CommentTok{\# Respondent degree distribution (projected network)}
\NormalTok{resp\_deg\_df }\OtherTok{\textless{}{-}} \FunctionTok{tibble}\NormalTok{(}\AttributeTok{degree =} \FunctionTok{degree}\NormalTok{(proj))}

\NormalTok{p2 }\OtherTok{\textless{}{-}} \FunctionTok{ggplot}\NormalTok{(resp\_deg\_df, }\FunctionTok{aes}\NormalTok{(}\AttributeTok{x =}\NormalTok{ degree)) }\SpecialCharTok{+}
  \FunctionTok{geom\_histogram}\NormalTok{(}
    \AttributeTok{binwidth =} \DecValTok{1}\NormalTok{, }\AttributeTok{fill =} \StringTok{"\#2A9D8F"}\NormalTok{, }\AttributeTok{alpha =} \FloatTok{0.8}\NormalTok{,}
    \AttributeTok{colour =} \StringTok{"white"}\NormalTok{, }\AttributeTok{boundary =} \SpecialCharTok{{-}}\FloatTok{0.5}
\NormalTok{  ) }\SpecialCharTok{+}
  \FunctionTok{geom\_vline}\NormalTok{(}
    \FunctionTok{aes}\NormalTok{(}\AttributeTok{xintercept =} \FunctionTok{mean}\NormalTok{(degree)),}
    \AttributeTok{linetype =} \StringTok{"dashed"}\NormalTok{, }\AttributeTok{colour =} \StringTok{"\#E63946"}
\NormalTok{  ) }\SpecialCharTok{+}
  \FunctionTok{labs}\NormalTok{(}
    \AttributeTok{x =} \StringTok{"Degree"}\NormalTok{, }\AttributeTok{y =} \StringTok{"Count"}\NormalTok{,}
    \AttributeTok{title =} \StringTok{"Respondent Degree Distribution}\SpecialCharTok{\textbackslash{}n}\StringTok{(Projected Network)"}
\NormalTok{  ) }\SpecialCharTok{+}
  \FunctionTok{theme}\NormalTok{(}\AttributeTok{plot.title =} \FunctionTok{element\_text}\NormalTok{(}\AttributeTok{face =} \StringTok{"bold"}\NormalTok{, }\AttributeTok{size =} \DecValTok{10}\NormalTok{))}

\CommentTok{\# Collaboration frequency}
\NormalTok{p3 }\OtherTok{\textless{}{-}}\NormalTok{ edges }\SpecialCharTok{|\textgreater{}}
  \FunctionTok{filter}\NormalTok{(freq\_clean }\SpecialCharTok{!=} \StringTok{"Unknown"}\NormalTok{) }\SpecialCharTok{|\textgreater{}}
  \FunctionTok{count}\NormalTok{(freq\_clean) }\SpecialCharTok{|\textgreater{}}
  \FunctionTok{ggplot}\NormalTok{(}\FunctionTok{aes}\NormalTok{(}\AttributeTok{x =}\NormalTok{ freq\_clean, }\AttributeTok{y =}\NormalTok{ n, }\AttributeTok{fill =}\NormalTok{ freq\_clean)) }\SpecialCharTok{+}
  \FunctionTok{geom\_col}\NormalTok{(}\AttributeTok{alpha =} \FloatTok{0.85}\NormalTok{, }\AttributeTok{width =} \FloatTok{0.7}\NormalTok{) }\SpecialCharTok{+}
  \FunctionTok{geom\_text}\NormalTok{(}\FunctionTok{aes}\NormalTok{(}\AttributeTok{label =}\NormalTok{ n), }\AttributeTok{vjust =} \SpecialCharTok{{-}}\FloatTok{0.4}\NormalTok{, }\AttributeTok{fontface =} \StringTok{"bold"}\NormalTok{, }\AttributeTok{size =} \FloatTok{3.5}\NormalTok{) }\SpecialCharTok{+}
  \FunctionTok{scale\_fill\_manual}\NormalTok{(}\AttributeTok{values =}\NormalTok{ freq\_palette) }\SpecialCharTok{+}
  \FunctionTok{labs}\NormalTok{(}\AttributeTok{x =} \ConstantTok{NULL}\NormalTok{, }\AttributeTok{y =} \StringTok{"Count"}\NormalTok{, }\AttributeTok{title =} \StringTok{"Collaboration Frequency"}\NormalTok{) }\SpecialCharTok{+}
  \FunctionTok{theme}\NormalTok{(}
    \AttributeTok{legend.position =} \StringTok{"none"}\NormalTok{,}
    \AttributeTok{plot.title =} \FunctionTok{element\_text}\NormalTok{(}\AttributeTok{face =} \StringTok{"bold"}\NormalTok{, }\AttributeTok{size =} \DecValTok{10}\NormalTok{)}
\NormalTok{  )}

\NormalTok{p1 }\SpecialCharTok{+}\NormalTok{ p2 }\SpecialCharTok{+}\NormalTok{ p3 }\SpecialCharTok{+}
  \FunctionTok{plot\_annotation}\NormalTok{(}
    \AttributeTok{title =} \StringTok{"Network Structure: Degree Distributions and Collaboration Patterns"}\NormalTok{,}
    \AttributeTok{theme =} \FunctionTok{theme}\NormalTok{(}\AttributeTok{plot.title =} \FunctionTok{element\_text}\NormalTok{(}\AttributeTok{face =} \StringTok{"bold"}\NormalTok{, }\AttributeTok{size =} \DecValTok{14}\NormalTok{))}
\NormalTok{  )}
\end{Highlighting}
\end{Shaded}

\begin{figure}[H]

{\centering \includegraphics[width=1\linewidth,height=\textheight,keepaspectratio]{sna_analysis-2_files/figure-pdf/degree-distributions-1.pdf}

}

\caption{Figure 5. Left: Organisation degree distribution. Centre:
Respondent degree distribution (projected). Right: Collaboration
frequency.}

\end{figure}%

\section{Discussion and Implications}\label{discussion-and-implications}

The Canadian disability rights community operates as a \textbf{highly
decentralized network}. Collaboration is predominantly local and
issue-specific rather than nationally integrated. A small number of
national organizations serve as critical connectors whose capacity to
convene and bridge across issue silos is disproportionately important to
overall network cohesion.

\subsection{Fragmentation as a structural
feature}\label{fragmentation-as-a-structural-feature}

The Canadian disability rights community, as captured by this survey,
operates as a highly decentralised network. Collaboration is
predominantly local and issue-specific rather than nationally
integrated. This is not necessarily a weakness --- it may reflect the
diversity of disability issues and the importance of place-based
advocacy --- but it does mean that information, resources, and strategic
coordination may not flow easily across the network.

\subsection{The role of national
organisations}\label{the-role-of-national-organisations}

A small number of national organisations (Disability Without Poverty,
CNIB, March of Dimes, Inclusion Canada) serve as critical connectors.
These organisations' capacity to convene, share information, and bridge
across issue silos is disproportionately important to overall network
cohesion. Strengthening these hub organisations --- or creating new ones
--- could substantially improve coordination across the community.

\subsection{Bridging actors}\label{bridging-actors}

Respondents with high betweenness centrality are structurally important:
they connect otherwise isolated parts of the network. Supporting and
amplifying these bridging actors --- many of whom are persons with
disabilities themselves --- could enhance cross-pollination of ideas and
strategies.

\subsection{Limitations}\label{limitations}

This analysis is based on a convenience sample of 111 respondents. The
free-text partner naming format means some organisations may be
inconsistently recorded. The projected one-mode network relies on the
assumption that shared organisational partnerships indicate a meaningful
indirect connection between respondents. Finally, the survey captures
reported collaboration, which may differ from actual interaction
patterns.

\begin{center}\rule{0.5\linewidth}{0.5pt}\end{center}

\emph{Report generated from the Enquête d'analyse des Réseaux Sociaux /
Social Network Analysis Survey. Analysis performed in R using igraph,
tidygraph, and ggraph. All respondent identifiers are de-identified
survey IDs.}




\end{document}
